\chapter{Metodologia}
\label{cap:metodologia}

\section{Desenvolvimento do Protótipo}

O desenvolvimento da Fase 1 do protótipo do sistema de treinamento em Realidade Virtual (RV) será orientado e validado por meio de consultas a especialistas em segurança do trabalho, incluindo engenheiros de segurança e instrutores de treinamentos de NR-18 \cite{brasil2025nr18} e NR-35 \cite{brasil2023nr35}, bem como representantes de empresas e entidades do setor da construção civil, como o \gls{SENAI} e o SINDUSCON-CE.

Essas consultas visam validar os principais elementos que devem ser priorizados nos cenários virtuais, considerando as principais causas de acidentes no setor, as dificuldades enfrentadas pelos instrutores na transmissão de conhecimentos e os temas de maior relevância para o público-alvo. O foco desta fase está delimitado aos riscos de quedas de altura e ao uso correto de Equipamentos de Proteção Individual (EPIs), por representarem as situações de maior impacto e maior potencial de aproveitamento das características imersivas da RV.

\section{Ferramentas Tecnológicas}

O sistema foi desenvolvido utilizando as seguintes tecnologias:

\begin{itemize}
    \item \textbf{Motor:} Unity 6 LTS (versão 6000.3.8f1)
    \item \textbf{Linguagem de Programação:} C\#
    \item \textbf{\gls{SDK} de RV:} Meta XR SDK (versão 201.0.0)
    \item \textbf{Dispositivo alvo:} Meta Quest 3, com rastreamento de mãos (\textit{hand tracking}) como modalidade primária de entrada.
\end{itemize}

\section{Estrutura dos Cenários Virtuais}

Os cenários de treinamento serão definidos a partir de consultas com especialistas em segurança do trabalho e representantes do setor da construção civil. Esta fase contempla dois cenários principais, alinhados às diretrizes da NR-18 \cite{brasil2025nr18} e NR-35 \cite{brasil2023nr35}:

\textbf{Cenário 1 -- Seleção de EPIs}

Cena introdutória na qual o usuário é recebido por um instrutor virtual e deve selecionar e equipar os EPIs adequados ao trabalho em altura que realizará em seguida. O grupo reúne cinco tarefas em modo sequencial, com limite de 270 segundos: equipar botina de segurança, luvas, óculos de proteção, capacete com jugular (exigido pela NR-35) e cinturão de segurança tipo paraquedista. A bancada apresenta, além dos itens corretos, distratores que forçam decisões fundamentadas na norma: capacete sem jugular, cinto abdominal (vedado pela NR-35 para trabalho em altura), luva de látex (inadequada para estruturas metálicas), calçado aberto (vedado pela NR-18) e colete refletivo (não pontuado). O objetivo é desenvolver o reconhecimento dos equipamentos obrigatórios antes da exposição ao cenário de risco.

\textbf{Cenário 2 -- Trabalho em Altura}

Cenário principal de inspeção de um andaime fachadeiro com irregularidades deliberadas. O grupo reúne quatro tarefas em modo de ordem livre, com limite de 300 segundos e dependência da conclusão do Cenário 1: conectar o talabarte a um ponto de ancoragem certificado (tipo A1, conforme a ABNT NBR 16325 \cite{nbr16325}), instalar guarda-corpo ausente (NR-18, item 18.15), instalar rodapé ausente (NR-18, item 18.15) e sinalizar a queda de uma rede de proteção. Caso o usuário inicie a atividade sem corrigir as irregularidades, o sistema dispara consequências --- como a queda de uma ferramenta pela borda desprotegida e uma simulação de queda por talabarte desconectado ---, cada uma registrada como violação de segurança com penalização na pontuação. \textit{Feedback} imediato é fornecido a cada ação.

O sistema prevê expansão modular para novos cenários em fases futuras, como queda de objetos, soterramentos e inspeção geral de canteiro, conforme validação com especialistas e parceiros institucionais.

\section{Amostragem e Grupos Experimentais}

A pesquisa contará com 30 participantes, divididos aleatoriamente em dois grupos:

\begin{itemize}
    \item \textbf{Grupo Experimental (RV):} Receberá o treinamento utilizando o sistema de Realidade Virtual.
    \item \textbf{Grupo Controle (Treinamento Tradicional):} Participará do treinamento por meio de vídeo instrucional com conteúdo equivalente ao apresentado no sistema de RV.
\end{itemize}

Os participantes poderão ser trabalhadores da construção civil ou estudantes da área. A alocação aleatória entre os grupos experimental e controle caracteriza o desenho como um \textbf{estudo piloto randomizado controlado}, de caráter exploratório e pioneiro no contexto brasileiro -- onde não foram identificados estudos controlados sobre RV aplicada à segurança do trabalho sob as normas regulamentadoras nacionais. Com $n = 30$, a análise prioriza a estimativa de \textbf{tamanhos de efeito} e respectivos intervalos de confiança, mais informativos que a significância isolada nesse tamanho amostral, fornecendo subsídios para o dimensionamento de pesquisas futuras com maior poder estatístico.

\section{Instrumentos de Avaliação}

Os instrumentos foram selecionados para cobrir as três hipóteses do estudo com o mínimo de sobrecarga para os participantes, garantindo viabilidade de aplicação em contexto de pesquisa piloto.

\subsection{Instrumentos aplicados a ambos os grupos}

\begin{itemize}
    \item \textbf{Teste de Conhecimento (pré e pós-teste):} Aplicado antes e imediatamente após o treinamento, mede o ganho de conhecimento sobre procedimentos de segurança em trabalhos em altura e uso correto de EPIs. As questões serão elaboradas com base nos conteúdos dos cenários desenvolvidos e validadas por especialistas em segurança do trabalho.
    \item \textbf{\textit{Reduced Instructional Materials Motivation Survey} (RIMMS):} Versão reduzida de 12 itens (três por subescala) do IMMS, baseada no modelo ARCS \cite{keller1987} e validada por \citeonline{loorbach2015}. Avalia motivação e engajamento nas dimensões de Atenção, Relevância, Confiança e Satisfação em escala Likert de cinco pontos. Será utilizada a versão traduzida e adaptada para o português brasileiro por \citeonline{cardosojunior2020}. Aplicado após o treinamento a ambos os grupos, permite comparação direta da experiência motivacional entre o método imersivo e o vídeo instrucional com menor carga de resposta em relação ao instrumento completo.
\end{itemize}

\subsection{Instrumentos aplicados exclusivamente ao grupo experimental}

\begin{itemize}
    \item \textbf{\textit{System Usability Scale} (SUS):} Instrumento padronizado para mensurar a usabilidade e aceitação do sistema pelos usuários \cite{brooke1996}. Será utilizada a versão brasileira validada por \citeonline{lourencovalentim2022}, que apresentou equivalência semântica, idiomática, conceitual e cultural com a versão original, com consistência interna de $\alpha = 0{,}76$.
    \item \textbf{\textit{iGroup Presence Questionnaire} (IPQ):} Avalia a sensação de presença e imersão no ambiente virtual \cite{schubert2001}. Será utilizada a versão portuguesa validada por \citeonline{vasconcelosraposo2016}, com adaptações lexicais mínimas para o português brasileiro.
\end{itemize}

\subsection{Teste de Retenção (condicional)}

Sujeito à viabilidade de cronograma, um teste de retenção de conhecimento equivalente ao pós-teste será aplicado entre 30 e 45 dias após o treinamento a ambos os grupos, com o objetivo de verificar a permanência do aprendizado ao longo do tempo. Caso inviável, esta etapa será declarada como limitação do estudo e recomendada para investigações futuras.

\section{Procedimento de Coleta}

A coleta de dados seguirá a seguinte sequência para ambos os grupos:

\begin{enumerate}
    \item Assinatura do \gls{TCLE}
    \item Aplicação do pré-teste de conhecimento
    \item Realização do treinamento (RV ou vídeo instrucional, conforme o grupo)
    \item Aplicação do pós-teste de conhecimento
    \item Aplicação do questionário de motivação e engajamento
    \item Aplicação do SUS e IPQ \textit{(somente grupo experimental)}
    \item Aplicação do teste de retenção \textit{(ambos os grupos, condicional)}
\end{enumerate}

\section{Avaliação Qualitativa}

Após a conclusão do treinamento, será conduzida uma análise qualitativa baseada em entrevistas semiestruturadas com os participantes do grupo experimental. Essas entrevistas terão como objetivos explorar a experiência subjetiva dos participantes, coletar percepções sobre o realismo e a utilidade do treinamento em RV, identificar pontos fortes e fracos da abordagem e obter sugestões de melhoria para fases futuras do sistema.
